\section*{历史注记}
\phantomsection
\addcontentsline{toc}{section}{历史注记}

（注意：罗马数字指本注记末尾的参考文献。）

对序列 \((u_{n})\) 用一个依赖于实参数 \(\lambda\) 且等于 \(u_{n}\) （当 \(\lambda = n\) 时）的积分的值进行``插值''的想法，可追溯到沃利斯（III, p.~55）。正是这一想法主要引导了欧拉，他在 1730 年（(I), v. XIV, p.~1-24）着手对阶乘序列进行插值。他首先注意到，\(n!\) 等于无穷乘积 \(\prod_{k=1}^{\infty}\left(\frac{k+1}{k}\right)^{n}\frac{k}{k+n}\) ，这个乘积对 n 的每个值（整数或非整数）都有定义，并且特别地，当 \(n = \frac{1}{2}\) 时由沃利斯公式可知它取值 \(\frac{1}{2}\sqrt{\pi}\) 。这一结果与沃利斯的结果之间的类比，促使他重新考察积分

\[
\int_ {0} ^ {1} x ^ {e} (1 - x) ^ {n} d x
\]

（n 为整数，e 任意），沃利斯已经考虑过它。欧拉利用二项展开得到它的值 \(\frac{n!}{(e+1)(e+2)\ldots(e+n+1)}\) ；一个变量替换又向他表明，\(n!\) 是积分 \(\int_{0}^{1}\left(\frac{1-x^{z}}{z}\right)^{n}dx\) 当 z 趋于 0 时的极限，由此得到``第二类欧拉积分''

\[
n! = \int_ {0} ^ {1} \left(\log \frac {1}{x}\right) ^ {n} d x;
\]

用同样的方法，并利用沃利斯公式，他得到公式 \(\int_0^1\sqrt{\log 1 / x} dx = \frac{1}{2}\sqrt{\pi}\) 。在他的后期著作中，欧拉经常回到这些积分；他就这样发现了余元公式（(I), t. XV, p.~82 与 t. XVII, p.~342）、公式 \(\mathbf{B}(p,q) = \Gamma (p)\Gamma (q) / \Gamma (p + q)\) （(I), t. XVII, p.~355）以及勒让德–高斯公式对应于 \(x = 1\) 的特殊情形（(I), t. XIX, p.~483）；所有这些当然都没有考虑收敛性问题。

高斯在关于超几何函数的研究中继续研究 \(\Gamma\) 函数，\(\Gamma\) 函数是超几何函数的一个极限情形（II）；正是在这项研究过程中他得到了一般的倍乘公式（勒让德稍早已对 p = 2 记下了它）。关于 \(\Gamma\) 的后期工作主要涉及把这个函数延拓到复数域。直到最近人们才认识到，对数凸这一性质在实数域上把 \(\Gamma(x)\) （相差一个常数因子）从泛函方程 \(f(x + 1) = x f(x)\) 的解中刻画出来（III）；阿廷（IV）指出了如何把 \(\Gamma(x)\) 的所有经典结果都简单地与这一性质联系起来。我们相当紧密地遵循了他的阐述。

